% ===== 原文 pp.28–29 =====
% 说明：本文件首行内容紧接 tr04a.tex 末尾的 "$\omega_{\what{\cT}_{g,n}}(-L_jt_j,\ )=$"，
%       该行是 p.27→p.28 跨页的行内公式，两文件之间不得插空行（否则公式会被断开成两段）。
% [原文 p.28]
$d(\half L_j^2)$（矩值 $\half L^2$ 确定了扭转的符号/定向）。向量场 $-L_jt_j$ 是 $(S^1)^n$ 作用在 $\what{\cT}_{g,n}$ 上的无穷小生成元，该作用由旋转边界点给出。辛形式 $\omega_{\what{\cT}_{g,n}}$ 是扭转不变的，因此我们已为辛约化做好准备，即准备好考察商 $\what{\cT}_{g,n}/(S^1)^n$。

% [Tutor Note: 这里给出的是"是 Hamilton 势"的验证链：由扭转–长度对偶 $\omega_{\what{\cT}_{g,n}}(-t_j,\ )=\dd L_j$，立刻得到 $\omega_{\what{\cT}_{g,n}}(-L_jt_j,\ )=L_j\dd L_j=\dd(\half L_j^2)$，所以 $\half L_j^2$ 确为 $-L_jt_j$ 的 Hamilton 势；这也解释了为什么 p.30 的 Duistermaat--Heckman 公式里出现的是系数 $L_j^2/2$ 而不是 $L_j$——矩映射正是 $\mu=\half L^2$（p.27 的 $\what{L}=(L_1^2/2,\dots,L_n^2/2)$）。]

\begin{figure}[htbp]
\centering
\includegraphics[width=0.75\textwidth]{fig10.png}
\figtitle{10}{$(S^1)^3$ 作用的正向旋转。}
\end{figure}

矩映射 $\mu:\what{\cT}_{g,n}\longrightarrow\R^n$ 的一个水平集，是 $\beta$ 长度取指定值的双曲曲面构成的轨迹。群 $(S^1)^n$ 通过旋转近乎紧的裤子作用在这些水平集上。水平集对该群的商自然地就是 $\cT_g(L)$——水平集规定了边界长度，而群作用去掉了关于这些点位置的信息。

\begin{prop}{11}
辛约化。若 $\what{\cT}_{g,n}/(S^1)^n\approx\cT_g(L)$，则
\[
\omega_{\what{\cT}_{g,n}}\big|_{\mu^{-1}(\what{L})/(S^1)^n}\approx\omega_{\cT_g(L)}.
\]
\end{prop}

\begin{proof}
形式 $\omega_{\what{\cT}_{g,n}}$ 由任何包含多重曲线 $\beta$ 的标准模型曲面的裤子分解的 $\dd\ell\wedge\dd\tau$ 公式给出。微分 $\dd L_j$ 在 $\mu$ 的水平集上消失，于是该公式约化为带边曲面的一个裤子分解所对应的和，这是一个不含边界参数的和——正是 $\omega_{\cT_g(L)}$ 的公式。
\end{proof}

% [Tutor Note: 命题 11 的证明只用到 $\dd\ell\wedge\dd\tau$ 公式的"逐项消失"：$\mu$ 的水平集上 $L_1,\dots,L_n$ 为常数，故 $n$ 个边界项 $\dd\ell_j\wedge\dd\tau(\beta_j)$ 中的 $\dd L_j=0$ 全部消失，剩下的恰是不含边界参数的那部分和。注意结论是"限制到 $\mu^{-1}(\what{L})/(S^1)^n$ 后形式相等"，与 p.30 用 Duistermaat--Heckman 得到的"上同调等价"是两种不同强度的陈述：前者在形式层面成立，后者只到上同调类（因为那里还要让 $L_j$ 变动）。]

\para{$S^1$ 主丛。}
我们回顾特征类的基本知识。

\begin{defn}{12}
设 $\pi:P\longrightarrow M$ 是光滑紧流形 $M$ 上的光滑圆周丛。若满足以下条件，则称该丛为 $S^1$ 主丛，
\begin{enumerate}[leftmargin=1.8em,itemsep=2pt,topsep=2pt]
\item $S^1$ 自由地作用在 $P$ 上，
\item $\pi(p_1)=\pi(p_2)$ 当且仅当存在 $s\in S^1$，使得 $p_1\cdot s=p_2$。
\end{enumerate}
$S^1$ 主丛的一个联络，是指由 $\mathbf{T}P$ 的切子空间构成的光滑分布 $\mathcal{H}$，使得
\begin{enumerate}[leftmargin=1.8em,itemsep=2pt,topsep=2pt]
\item 对每个 $p\in P$，有 $\mathbf{T}_pP=\mathcal{H}_p\oplus\ker\pi_*|_p$，
% [原文 p.29]
\item $s^*\mathcal{H}_p=\mathcal{H}_{p\cdot s}$。
\end{enumerate}
% 注：以下一段在原文 p.29 仍用斜体排印，属定义 12 的正文（联络与曲率形式的说明），故留在 defn 环境内。
% [Tutor Note: 定义 12 第 2 条里的群作用方向值得留意：条件是 $\pi(p_1)=\pi(p_2)\iff\exists s\in S^1,\ p_1\cdot s=p_2$，即 $S^1$ 从\emph{右}作用、并把同一纤维上的点传递地彼此搬动；配合第 1 条"作用自由"，即得 $S^1$ 在每条纤维上自由且传递（纤维是 $S^1$-torsor）。这也解释了下一行 $s^*\mathcal{H}_p=\mathcal{H}_{p\cdot s}$ 的下标写成 $p\cdot s$（而不是 $s\cdot p$）：联络的 $S^1$ 不变性是就右作用写的。两种写法只差取逆 $s\mapsto s^{-1}$，但全篇必须统一，否则 $A(\dot{s})=1$ 与 $c_1$ 的符号会跟着翻。]

联络唯一地由 $\mathcal{H}=\ker A$ 给出，其中 $A$ 是 $P$ 上的 1-形式，只要 $A$ 是 $S^1$ 不变的且 $A(\dot{s})=1$。$S^1$ 不变的内积 $\langle\,,\rangle$ 通过 $A(v)=\langle v,\dot{s}\rangle/\langle\dot{s},\dot{s}\rangle$ 提供不变 1-形式的一个例子。联络在 $P$ 上的曲率 2-形式为 $\Phi(v,w)=\dd A(\mathrm{hor}\,v,\mathrm{hor}\,w)$，其中 hor 表示 $\mathbf{T}P$ 到 $\mathcal{H}$ 的\emph{水平}投影。
\end{defn}

\begin{thm}{13} [MS74]。存在 $M$ 上唯一的闭 2-形式 $\Omega$，使得 $\Phi=\pi^*\Omega$。$\Omega$ 的上同调类与 $P$ 的 $S^1$ 主联络的选取无关，且第一 Chern 类为 $c_1(P)=[\Omega]\in H^2(M,\Z)$。\end{thm}

% [Tutor Note: 定理 13 中"$\Omega$ 唯一"要理解为：$\Phi$ 是 $P$ 上的形式，而 $\pi^*:H^\bullet(M)\to H^\bullet(P)$ 单射（纤维可缩），故一旦给定联络、便存在唯一的 $\Omega$ 使 $\Phi=\pi^*\Omega$——唯一性是对固定的 $\Phi$ 而言，不是对所有联络而言；跨联络的不变性只对上同调类 $[\Omega]$ 成立，这正是定理后半句的内容。另一处需注意记号约定：$c_1(P)=[\Omega]$ 成立的前提是把 $\Omega$ 按 $1/2\pi$ 归一（等价于把 $\dot{s}$ 取为单位速度、$A(\dot{s})=1$）；若采用 $[\,\Omega/2\pi\,]\in H^2(M,\Z)$ 的写法，与本文相差一个 $2\pi$ 因子，引用到体积公式（定理 10）时须核对同一约定。]

如上，曲面 $R$ 的边界 $\beta_j$ 上的变动点定义了 $\cT_g(L)$ 上的 $S^1$ 主丛 $\what{\beta}_j$；$S^1$ 通过让该点移动、且曲面内部位于左侧来作用。为丛选取一个联络，就给出第一 Chern 类 $c_1(\what{\beta}_j)$。

\para{应用 Duistermaat--Heckman 定理。}
我们把 $\what{\cT}_{g,n}$ 的定义扩充到包含 $L=0$ 的情形；$R$ 的测地边界可以换成穿刺点。当边界长度趋于零时双曲结构收敛；特别地，领口区域收敛到尖点区域。$\what{\cT}_{g,n}$ 的扩充，由"用领口/尖点区域边界上的点来参数化边界点"给出。

我们回顾领口与尖点的基本知识。对长度为 $\ell_\alpha$ 的测地线 $\alpha$，上半平面 $\mathbb{H}$ 中的标准领口是 $\mathcal{C}(\ell_\alpha)=\{d(z,i\R_+)\le w(\alpha)\}$，其中半宽 $w(\alpha)$ 由 $\sinh w(\alpha)\sinh\ell_\alpha/2=1$ 给出。商柱面 $\{d(z,i\R_+)\le w(\alpha)\}/\langle z\mapsto e^{\ell_\alpha}z\rangle$ 嵌入 $R$，给出该测地线的一个领口邻域。对尖点而言，$\mathbb{H}$ 中的标准尖点是 $\mathcal{C}_\infty=\{\Im z\ge 1/2\}$。商柱面 $\{\Im z\ge 1/2\}/\langle z\mapsto z+1\rangle$ 嵌入 $R$，给出一个尖点区域。当 $\ell_\alpha$ 有界时，领口的边界、以及尖点区域的边界，长度约为 2。领口与尖点区域都被与边界正交的测地线叶状分层（foliated）。当测地长度趋于零时，半领口邻域 Gromov--Hausdorff 收敛到尖点区域（收敛在有界距离的边界邻域上是一致的）；边界以及与边界正交的测地线都收敛。

% [Tutor Note: "长度约为 2"可以定量核对：标准领口中到轴距离为 $w$ 的等距曲线长度为 $\ell_\alpha\cosh w$，由 $\sinh w(\alpha)\sinh(\ell_\alpha/2)=1$ 解出 $\cosh w=\coth(\ell_\alpha/2)$，故领口边界长度为 $\ell_\alpha\coth(\ell_\alpha/2)$；当 $\ell_\alpha\to0$ 时它 $\to2$（把 $\coth(\ell/2)\sim2/\ell$ 代入即得），而尖点区域边界 $\Im z=1/2$ 在度量 $|\dd z|/\Im z$ 下的长度恰为 $1/(1/2)=2$——两个极限值完全吻合。这既解释了"约为 2"，也定量说明了原文"领口区域收敛到尖点区域"。]

\begin{figure}[htbp]
\centering
\includegraphics[width=0.7\textwidth]{fig11.png}
\figtitle{11}{沿测地线投影到领口（collar）与尖点（cusp）区域边界。}
\end{figure}

与领口边界正交的测地线，给出从核心测地线到每条领口边界的投影。该投影用来指出：
